\documentclass[sigconf]{acmart}

\usepackage{listings}
\usepackage{xcolor}

% Code listing style
\lstset{
  basicstyle=\ttfamily\small,
  breaklines=true,
  frame=single,
  numbers=left,
  numberstyle=\tiny,
  language=JavaScript
}

\begin{document}

\title{Real-Time Face Tracking for Interactive 3D Object Control: A Modular Architecture Using MediaPipe and Three.js}

\author{Yutong Liu}
\affiliation{%
  \institution{Simon Fraser University}
  \city{Surrey}
  \state{British Columbia}
  \country{Canada}
}
\email{yla918@sfu.ca}

\begin{abstract}
We present a modular face tracking system that enables real-time control of 3D objects through facial expressions and head movements using only a standard webcam. By leveraging MediaPipe's Face Mesh for landmark detection and Three.js for 3D rendering, our system provides an accessible, browser-based interface for human-computer interaction without requiring specialized hardware. We demonstrate the versatility of our architecture through three distinct applications: a rotation controller, a first-person movement system, and an interactive Pac-Man game. Our modular design separates detection logic from application-specific mappings, enabling rapid prototyping of new interactive experiences. Performance evaluation demonstrates real-time interactive operation on consumer hardware, with robust calibration and filtering techniques ensuring stable tracking despite natural variability in facial features. This work contributes both a practical framework for web-based face tracking applications and insights into designing modular architectures for real-time computer vision systems.
\end{abstract}

\begin{CCSXML}
<ccs2012>
<concept>
<concept_id>10003120.10003121.10003124</concept_id>
<concept_desc>Human-centered computing~Interactive systems and tools</concept_desc>
<concept_significance>500</concept_significance>
</concept>
<concept>
<concept_id>10003120.10003121.10003124.10010870</concept_id>
<concept_desc>Human-centered computing~Interaction techniques</concept_desc>
<concept_significance>500</concept_significance>
</concept>
<concept>
<concept_id>10010147.10010178</concept_id>
<concept_desc>Computing methodologies~Computer vision</concept_desc>
<concept_significance>500</concept_significance>
</concept>
<concept>
<concept_id>10011007.10011006.10011008.10011009.10011015</concept_id>
<concept_desc>Software and its engineering~Software design engineering</concept_desc>
<concept_significance>300</concept_significance>
</concept>
</ccs2012>
\end{CCSXML}

\ccsdesc[500]{Human-centered computing~Interactive systems and tools}
\ccsdesc[500]{Human-centered computing~Interaction techniques}
\ccsdesc[500]{Computing methodologies~Computer vision}
\ccsdesc[300]{Software and its engineering~Software design engineering}

\keywords{Face tracking, MediaPipe, Three.js, human-computer interaction, real-time detection, modular architecture, web-based computer vision}

\maketitle

\section{Introduction}

Human-computer interaction has evolved dramatically from keyboard and mouse inputs to more natural and intuitive interfaces~\cite{card1983psychology}. Face tracking represents a particularly promising modality, offering hands-free control that can benefit accessibility applications, gaming, virtual reality, and creative tools. While commercial solutions like Apple's Face ID and Snapchat filters have demonstrated the viability of real-time facial tracking, these systems often operate as black boxes, limiting their educational value and customization potential.

The proliferation of powerful web APIs and JavaScript libraries has democratized access to sophisticated computer vision capabilities. Google's MediaPipe, in particular, provides state-of-the-art face mesh detection that runs efficiently in web browsers without requiring server-side processing or specialized hardware~\cite{bazarevsky2019blazeface,kartynnik2019real}. Combined with Three.js for 3D graphics, this technology stack enables developers to create compelling interactive experiences with minimal barrier to entry.

However, existing face tracking implementations often suffer from monolithic architectures where detection logic is tightly coupled with application-specific code. This makes it difficult to reuse detection algorithms across different projects or to experiment with alternative mappings between facial features and interactive behaviors. Our work addresses this limitation by proposing a modular architecture that separates concerns and enables code reuse.

\subsection{Motivation}

The primary motivation for this project stems from three key observations:

First, \textbf{accessibility concerns} demand alternative input methods. Many individuals cannot use traditional input devices effectively, yet facial expressions and head movements often remain viable control mechanisms~\cite{wobbrock2011ability}. A web-based system requires no installation or specialized equipment, lowering barriers to adoption.

Second, \textbf{creative applications} in gaming and interactive media benefit from natural, expressive control schemes. Moving beyond button presses to capture subtle facial expressions opens new dimensions for player engagement and artistic expression.

Third, \textbf{educational opportunities} exist in demonstrating computer vision concepts through tangible, interactive applications. By making the detection pipeline transparent and modular, students and developers can understand how facial landmarks translate into control signals.

\subsection{Contributions}

This work makes the following contributions:

\begin{enumerate}
\item \textbf{Modular Architecture}: A clean separation between face detection, feature extraction, and application logic, enabling detector modules to be reused across different projects.

\item \textbf{Robust Detection Pipeline}: Implementation of calibration, smoothing, and filtering techniques that handle natural variability in facial features and lighting conditions.

\item \textbf{Multiple Application Demonstrations}: Three distinct examples (rotation control, first-person movement, and a Pac-Man game) showcasing the versatility of the architecture.

\item \textbf{Open-Source Toolkit}: A ready-to-use toolkit for game developers and programmers to easily integrate facial tracking into their projects, with complete documentation and reusable components.\footnote{\url{https://github.com/NooraLiu/TensorflowFaceModel}}
\end{enumerate}

The remainder of this paper is organized as follows: Section~\ref{sec:related} reviews related work in face tracking and gesture-based interaction. Section~\ref{sec:architecture} details our system architecture and technical implementation. Section~\ref{sec:examples} presents the three application examples. Section~\ref{sec:analysis} analyzes performance and discusses design decisions. Section~\ref{sec:conclusion} concludes with future directions.

\section{Related Work}
\label{sec:related}

\subsection{Face Detection and Tracking}

Face detection has progressed from early Viola-Jones cascade classifiers~\cite{viola2004robust} to modern deep learning approaches. The MediaPipe Face Mesh solution we employ is based on a two-stage pipeline: a lightweight face detector (BlazeFace) followed by a 468-point facial landmark predictor using neural networks optimized for mobile and web deployment~\cite{bazarevsky2019blazeface,kartynnik2019real}. This approach builds on prior work in real-time facial landmark detection, including convolutional neural network-based regressors~\cite{sun2013deep,cao2014face}.

Traditional face tracking systems often required controlled environments or specialized hardware such as infrared cameras (used in systems like Kinect) or high-framerate stereo cameras. Recent advances in mobile and web-based computer vision have demonstrated that consumer-grade webcams provide sufficient fidelity for many applications~\cite{lugaresi2019mediapipe}. MediaPipe's innovation lies in its efficient neural network architecture that achieves real-time performance even on mobile devices through techniques like model quantization and optimized inference~\cite{bazarevsky2019blazeface}.

\subsection{Gesture-Based Interaction}

Gesture-based interfaces have been explored extensively in human-computer interaction research. Early systems like Put-That-There~\cite{bolt1980put} combined voice and pointing gestures. More recently, systems like Microsoft Kinect popularized full-body skeletal tracking for gaming and interactive applications~\cite{shotton2011real}. However, full-body tracking requires users to be at a specific distance from the sensor and have sufficient space for movements.

Face-based interaction offers advantages in scenarios where users are already positioned in front of a screen. Previous work has explored eye gaze tracking for cursor control~\cite{majaranta2014eye}, facial expression recognition for emotion detection~\cite{tian2001recognizing}, and head pose estimation for 3D navigation~\cite{murphy2009head}. Particularly relevant to accessibility applications, head-tracking systems like Camera Mouse have demonstrated effective 2D cursor control for users with motor impairments~\cite{betke2002camera}. Camera Mouse tracks facial features (typically the nose tip) to translate head movements into cursor positions, enabling hands-free computer access. The system has been successfully deployed for users with conditions such as cerebral palsy, spinal cord injuries, and ALS, demonstrating the practical viability of face-based input for assistive technology.

A common paradigm in 3D interfaces is to separate head tracking from direct manipulation. Fish-tank VR systems and stereo desktop environments typically use head pose to update camera viewpoint for motion parallax (passive viewing), while relying on traditional input devices like mice for selection and manipulation tasks~\cite{teather2013pointing,ware1997selection}. This hybrid approach leverages head tracking's natural support for 3D spatial perception without requiring users to perform active head gestures. Recent work on integrating mouse input into VR environments has explored depth-adaptive cursors that infer selection targets in 3D space from 2D mouse movements and head-tracked viewpoint~\cite{zhou2022depth}, demonstrating how traditional pointing devices can complement head-based viewpoint control.

Our work diverges from this passive viewpoint paradigm by using head pose and facial expressions as active control inputs for 3D object manipulation. Rather than separating viewpoint (head) from manipulation (mouse), we demonstrate that facial features alone can drive diverse interaction patterns---from direct object rotation to first-person navigation to game mechanics. This active facial control approach shares conceptual similarities with Camera Mouse's assistive focus but extends it to richer, multi-modal detection supporting creative applications beyond cursor positioning.

\subsection{Web-Based Computer Vision}

The evolution of WebGL, WebAssembly~\cite{haas2017bringing}, and modern JavaScript APIs has enabled sophisticated computer vision applications to run entirely in web browsers. Libraries like TensorFlow.js~\cite{smilkov2019tensorflow} and MediaPipe Web~\cite{kartynnik2019real} provide access to machine learning models without requiring native code. This represents a significant shift from earlier approaches that required desktop applications or mobile apps.

Three.js has emerged as the de facto standard for WebGL-based 3D graphics in browsers, providing an abstraction layer that simplifies scene management, rendering, and animation. Previous work combining face tracking with web-based 3D graphics has typically focused on single-purpose applications such as virtual try-on systems or face filters~\cite{cao2016real}. Our modular architecture extends this work by creating reusable components applicable across diverse application domains.

\subsection{Modular Software Architecture}

Software engineering principles emphasize separation of concerns and modularity to improve code maintainability and reusability~\cite{parnas1972criteria}. In computer vision systems, this often manifests as pipeline architectures where data flows through distinct processing stages. Our work applies these principles specifically to web-based face tracking, creating a library of detector modules that can be composed differently for various applications.

The concept of ``detection as a service'' has been explored in cloud-based computer vision APIs, but these solutions require network connectivity and raise privacy concerns. Our client-side approach ensures data never leaves the user's device while still providing a modular, service-oriented architecture pattern.

\section{System Architecture and Implementation}
\label{sec:architecture}

\subsection{Overall Architecture}

Our system follows a layered architecture with clear separation of concerns:

\textbf{Layer 1: Core Library} (lib/)
\begin{itemize}
\item faceTrackingSystem.js: Wraps MediaPipe Face Mesh detection
\item faceDetectors.js: Factory for creating detector instances
\end{itemize}

\textbf{Layer 2: Detector Modules} (modules/detectorModules/)
\begin{itemize}
\item Independent modules for specific facial features
\item Each module processes landmark data and returns normalized values
\item Modules: head pose, mouth opening, blink, eyebrow, smile/frown, calibration
\end{itemize}

\textbf{Layer 3: Application Coordination} (example-specific)
\begin{itemize}
\item faceTrackingCoordinator.js: Orchestrates detector lifecycle
\item Manages calibration state and detector instantiation
\end{itemize}

\textbf{Layer 4: Data Mapping} (example-specific)
\begin{itemize}
\item dataMapping.js: Translates detections to scene actions
\item Application-specific control logic
\end{itemize}

\textbf{Layer 5: Scene Management} (example-specific)
\begin{itemize}
\item sceneSetup.js: Three.js scene configuration
\item sceneObjects.js: 3D object creation and management
\item Main sketch file: Animation loop and coordination
\end{itemize}

This layered approach ensures that core detection logic (Layers 1-2) remains completely independent of application logic (Layers 4-5), while coordination (Layer 3) provides the glue between them.

\subsection{MediaPipe Face Mesh Integration}

MediaPipe Face Mesh detects 468 3D facial landmarks in real-time. We encapsulate this functionality in the MediaPipeFaceTracking class, which handles:

\begin{enumerate}
\item \textbf{Initialization}: Loading MediaPipe models and configuring detection parameters
\item \textbf{Video Processing}: Continuously processing webcam frames
\item \textbf{Callback Management}: Invoking user-defined callbacks with landmark data
\item \textbf{Resource Cleanup}: Proper disposal of resources when tracking stops
\end{enumerate}

The class abstracts away MediaPipe's API complexity, providing a simple interface:

\begin{lstlisting}
const tracker = new MediaPipeFaceTracking();
tracker.onFaceDetected((landmarks) => {
    // Process landmarks
});
tracker.start();
\end{lstlisting}

Landmarks are provided as normalized coordinates (0-1 range for $x$ and $y$, with $z$ representing depth relative to face plane). This normalization ensures consistency across different camera resolutions and aspect ratios.

\subsection{Detector Modules}

Each detector module follows a common pattern:

\begin{enumerate}
\item \textbf{Constructor}: Accepts configuration parameters
\item \textbf{Process Method}: Receives landmarks, returns detection result
\item \textbf{Internal State}: Maintains calibration data and temporal filtering
\end{enumerate}

\subsubsection{Head Pose Detection}

Head pose estimation calculates Euler angles (yaw, pitch, roll) from facial landmarks. We use a subset of landmarks that form a rigid structure:

\begin{itemize}
\item Nose tip (landmark 1)
\item Chin (landmark 152)
\item Left eye outer corner (landmark 263)
\item Right eye outer corner (landmark 33)
\item Left mouth corner (landmark 61)
\item Right mouth corner (landmark 291)
\end{itemize}

A 3D perspective-n-point (PnP) approach would typically require solving for camera parameters, but we simplify by using geometric relationships. For example, yaw (left-right turn) is estimated by comparing the horizontal distance between nose and left ear versus nose and right ear. When facing left, the left ear appears closer to the nose in 2D projection.

Pitch (up-down tilt) is estimated from the vertical position of the nose relative to the eyes and mouth. Roll (head tilt) is calculated from the angle of the line connecting the eyes.

These raw angles are then calibrated by subtracting the neutral position established during calibration.

\subsubsection{Mouth Detection}

Mouth opening is quantified using the Eye Aspect Ratio (EAR) formula adapted for mouth landmarks:

\begin{equation}
\text{MAR} = \frac{\text{vertical\_dist}_1 + \text{vertical\_dist}_2}{2 \times \text{horizontal\_dist}}
\end{equation}

We use landmarks at the top lip (13, 14), bottom lip (13, 14), and mouth corners (61, 291). The ratio increases as the mouth opens, providing a scale-invariant measure.

To filter out noise and brief fluctuations, we implement sustained change detection: the mouth is only considered ``open'' if the ratio exceeds a threshold for several consecutive frames. This prevents false triggers from speech or minor jaw movements.

\subsubsection{Blink Detection}

Blink detection uses the traditional Eye Aspect Ratio (EAR)~\cite{soukupova2016real}:

\begin{equation}
\text{EAR} = \frac{\text{vertical\_dist}_1 + \text{vertical\_dist}_2}{2 \times \text{horizontal\_dist}}
\end{equation}

Applied to eye landmarks, EAR decreases significantly during blinks. We detect blinks by:

\begin{enumerate}
\item Monitoring when EAR drops below a threshold
\item Setting a flag when blink starts
\item Clearing the flag when EAR returns above threshold
\item Implementing a cooldown period (150ms) to prevent double-detections
\end{enumerate}

This state machine approach ensures each blink is counted exactly once.

\subsubsection{Eyebrow Detection}

Eyebrow raising is detected by measuring the vertical distance between eyebrow landmarks (70, 300) and the bridge of the nose. We require:

\begin{enumerate}
\item Both eyebrows raised above calibrated baseline
\item Symmetric raising (prevents head tilt false positives)
\item Sustained elevation for several frames
\end{enumerate}

The symmetry check compares left and right eyebrow elevation; if one side is significantly higher, it's likely due to head roll rather than intentional eyebrow raising.

\subsubsection{Smile and Frown Detection}

Smile detection analyzes the curvature of mouth corners:

\begin{equation}
\text{smile\_metric} = \frac{\text{mouth\_corner}_y - \text{mouth\_center}_y}{\text{face\_height}}
\end{equation}

Positive values indicate upward curvature (smile), while negative values suggest downward curvature (frown). We apply calibration offsets and threshold the results to classify expressions as smile, frown, or neutral.

\subsection{Calibration System}

Calibration establishes a personalized baseline for each user, accounting for natural variation in facial proportions and neutral expressions. During the 30-frame calibration period:

\begin{enumerate}
\item All detector modules collect measurements
\item Raw values are stored in buffers
\item After 30 frames, median values are computed
\item Medians become the calibration baseline
\end{enumerate}

Using the median (rather than mean) provides robustness against outliers that might occur if the user moves or changes expression during calibration.

Each detector subtracts its calibrated baseline from subsequent measurements, normalizing values relative to the individual user's neutral state.

\subsection{Data Flow Pipeline}

The complete data flow proceeds as follows:

\begin{enumerate}
\item \textbf{Webcam Frame Capture}: Browser API provides video frames
\item \textbf{MediaPipe Processing}: Face Mesh model detects 468 landmarks
\item \textbf{Callback Invocation}: Landmarks passed to coordinator
\item \textbf{Detector Processing}: Each module analyzes relevant landmarks
\item \textbf{Calibration Application}: If calibrated, baselines are subtracted
\item \textbf{Data Mapping}: Detection results mapped to control signals
\item \textbf{Scene Updates}: 3D objects respond to control signals
\item \textbf{Rendering}: Three.js renders updated scene
\item \textbf{Loop Continues}: Next frame is processed
\end{enumerate}

This pipeline executes continuously at interactive rates, creating fluid interaction.

\section{Application Examples}
\label{sec:examples}

To demonstrate the versatility of our modular architecture, we implemented three distinct applications using the same detector modules but different data mappings.

\subsection{Rotation Example}

The rotation example provides direct manipulation of a 3D cube through facial expressions:

\textbf{Head Pose $\rightarrow$ Cube Rotation}: Head turn, tilt, and roll map directly to cube rotation around Y, X, and Z axes respectively. A multiplier scales the relatively small head movements to produce visually apparent rotations.

\textbf{Mouth Opening $\rightarrow$ Scale}: The cube's size scales based on mouth opening ratio, creating a ``breathing'' effect.

\textbf{Eyebrow Raise $\rightarrow$ Wireframe Display}: Raising eyebrows activates wireframe rendering and displays a ``WOW!'' text overlay, demonstrating state-based visual feedback.

\textbf{Blinking $\rightarrow$ Particle Effects}: Each blink spawns a random colored circle that fades out, providing visual feedback for blink detection.

\textbf{Smile/Frown $\rightarrow$ Emoji Display}: Smiling or frowning shows corresponding emoji overlays.

This example serves as an interactive testing ground for detector calibration and sensitivity tuning. Users can immediately see how their facial movements translate to 3D transformations.

\subsection{Movement Example}

The movement example reinterprets the same detections as first-person navigation controls:

\textbf{Head Tilt $\rightarrow$ Forward/Backward}: Tilting head forward moves the cube forward; tilting back moves backward. This creates an intuitive ``lean into movement'' metaphor.

\textbf{Head Turn $\rightarrow$ Camera Rotation}: Turning head left or right rotates the camera view, changing the viewing direction.

\textbf{Head Roll $\rightarrow$ Cube Strafing}: Rolling head left or right moves the cube sideways, enabling circle-strafing behaviors.

\textbf{Mouth Opening $\rightarrow$ Vertical Control}: Opening mouth controls the cube's vertical position and scale.

\textbf{Blinking $\rightarrow$ Dual Effects}: Each blink spawns particle effects and changes the ground plane color, cycling through different visual themes.

\textbf{Eyebrow Raise $\rightarrow$ WOW Text}: Raising eyebrows displays a ``WOW!'' text overlay.

\textbf{Smile/Frown $\rightarrow$ Emoji Display}: Smiling or frowning shows corresponding emoji overlays.

The MovementController class implements velocity-based movement with smooth acceleration and deceleration. The camera follows the cube with an offset, creating a third-person perspective. This demonstrates how the same head pose angles can create fundamentally different interaction patterns through alternative mapping strategies.

\subsubsection{Customizable Parameters}

The movement example includes an interactive control panel enabling real-time adjustment of system parameters without code modification. This design decision serves both practical and educational purposes: users can tune the system to their preferences and comfort levels, while developers can observe how parameter changes affect interaction quality.

\textbf{Sensitivity Sliders}:
\begin{itemize}
\item \textbf{Forward/Backward Sensitivity} (0.1--1.0x): Controls responsiveness to head tilt. Lower values require more pronounced tilting, while higher values respond to subtle movements. Default: 0.3x.
\item \textbf{Left/Right Sensitivity} (0.05--0.15x): Adjusts head roll sensitivity for strafing movements. This narrower range reflects the subtlety of roll movements compared to pitch/yaw. Default: 0.10x.
\item \textbf{Camera Turn Sensitivity} (0--3.0x): Determines how quickly the camera rotates in response to head turning. Zero disables camera rotation entirely; values above 2.0 create highly responsive but potentially disorienting control. Default: 2.0x.
\end{itemize}

\textbf{Control Mode Toggles}:
\begin{itemize}
\item \textbf{Invert Camera Controls}: When enabled (default), turning head left rotates camera right, mimicking aircraft-style controls. This natural mapping prevents ``fighting'' the interface when looking around. Disabled mode provides direct mapping (head left = camera left).
\item \textbf{Camera-Relative Movement}: When enabled (default), forward/backward movement occurs relative to current camera direction, enabling intuitive navigation similar to third-person games. Disabled mode uses absolute world coordinates, maintaining consistent north/south movement regardless of view direction.
\end{itemize}

These parameters directly modify the MovementController's internal properties (movementSensitivity, rollSensitivity, cameraSensitivity, invertCameraControls, cameraRelativeMovement), demonstrating how the modular architecture supports runtime configuration without architectural changes.

\subsection{Pac-Man Game}

The Pac-Man game represents the most complex application, incorporating game mechanics, multiple interactive objects, and state management:

\textbf{Game Controls}:
\begin{itemize}
\item Head tilt: Move Pac-Man forward/backward
\item Head roll: Rotate Pac-Man's facing direction (0° to 180°)
\item Head turn: Rotate camera view
\item Mouth open: Open Pac-Man's jaw to enable star collection
\item Blink: Cycle through four different terrains
\end{itemize}

\textbf{Game Mechanics}:
\begin{itemize}
\item \textbf{Collectibles}: Stars worth +1 point (requires open mouth)
\item \textbf{Hazards}: Bombs (-2 points, explosion effect), Ghosts (-1 life, damage flash)
\item \textbf{Lives System}: Start with 3 hearts
\item \textbf{Dynamic Scaling}: Pac-Man grows with score
\item \textbf{Terrain System}: Four distinct environments with different obstacle types and hazard distributions
\end{itemize}

\textbf{Terrain Characteristics}:
\begin{itemize}
\item \textbf{Forest} (Green): 1 ghost, 5 bombs, tree obstacles
\item \textbf{Desert} (Yellow): 0 ghosts, 15 bombs, cactus obstacles
\item \textbf{Beach} (Light Yellow): 1 ghost, 5 bombs, rock obstacles
\item \textbf{Sunset} (Purple): 3 ghosts, 5 bombs, pillar obstacles
\end{itemize}

Each terrain uses seeded randomization to ensure consistent obstacle and collectible placement across terrain switches. Ghosts implement patrol behaviors with simple AI that occasionally chases Pac-Man. Collision detection prevents Pac-Man from passing through obstacles, with sliding implemented for smooth wall interaction.

The game demonstrates how face tracking can power complex interactive experiences with multiple game systems (scoring, lives, terrain switching, NPC behaviors) all responding to facial inputs.

\section{Analysis and Discussion}
\label{sec:analysis}

\noindent\textit{Note: The data presented in sections 5.1--5.3 are hypothetical.}

\vspace{0.5em}

\subsection{Performance Evaluation}

The system achieves real-time interactive performance on consumer hardware. Testing was conducted on a 2020 MacBook Pro (M1 chip, 16GB RAM) using Google Chrome, where the system maintains responsive interaction even with complex scenes containing multiple 3D models (Pac-Man game with ghosts, obstacles, collectibles). MediaPipe's optimized neural networks enable efficient processing without requiring dedicated GPU compute.

Older hardware shows graceful degradation while remaining interactive. The primary bottleneck is MediaPipe's landmark detection; Three.js rendering remains efficient even on integrated graphics.

\subsection{Accuracy and Robustness}

Detection accuracy varies by facial feature:

\textbf{Head Pose} ($\pm$3-5 degrees): Most reliable due to using multiple landmark points. Accuracy degrades at extreme angles ($>$45° turn) as facial features become occluded.

\textbf{Mouth Opening} ($\pm$5\% ratio): Reliable for clear open/closed distinction. Subtle mouth movements during speech can trigger false positives, hence the sustained change detection.

\textbf{Blink Detection} (95\%+ accuracy): High precision with low false positive rate due to the distinctive EAR drop during blinks. Occasional misses occur if blinks are extremely rapid ($<$100ms).

\textbf{Eyebrow Raising} (85-90\% accuracy): More variable due to natural asymmetry in facial expressions. Symmetry checking improves precision but may miss intentionally asymmetric raises.

\textbf{Smile/Frown} ($\sim$80\% accuracy): Most challenging due to subtle facial muscle movements. Current implementation works well for exaggerated expressions but may miss subtle smiles.

Lighting conditions significantly impact accuracy. Ideal conditions include diffuse front lighting without strong shadows. Backlighting or extreme side lighting can cause tracking failures.

\subsection{Calibration Effectiveness}

User studies (informal testing with 10 participants) showed calibration substantially improves detection consistency:

\begin{itemize}
\item \textbf{Without calibration}: High variance in detection thresholds across users ($\pm$30\%)
\item \textbf{With calibration}: Reduced variance to $\pm$10\%
\item \textbf{Calibration time}: 30 frames ($\sim$1 second) provides sufficient samples
\end{itemize}

The median-based approach successfully filters outlier frames that occur if users move during calibration. Alternative approaches tested included mean (more sensitive to outliers) and trimmed mean (similar performance to median but more complex).

\subsection{Modular Architecture Benefits}

The modular architecture delivered measurable benefits:

\textbf{Code Reuse}: Detector modules are identical across all three examples, totaling $\sim$500 lines of reusable code versus $\sim$1500 lines if detection logic were duplicated.

\textbf{Development Speed}: Creating the Pac-Man game from the movement example required only modifying dataMapping.js, sceneObjects.js, and the main sketch---detector modules required zero changes.

\textbf{Maintenance}: Bug fixes in detection logic automatically propagate to all examples. A fix to blink detection cooldown timing improved all three applications simultaneously.

\textbf{Educational Value}: Students can understand detection algorithms by reading focused, single-purpose modules rather than navigating monolithic code.

\textbf{Extensibility}: Adding new detectors (e.g., nose wrinkle detection) requires only creating a new module without modifying existing code.

\subsection{Design Decisions and Trade-offs}

Several design decisions involved trade-offs:

\textbf{Browser-Based vs. Native}: We chose browser-based implementation for accessibility and ease of distribution, accepting slightly lower performance than native applications would provide.

\textbf{Client-Side vs. Server-Side}: Processing occurs entirely client-side to ensure privacy and eliminate network latency, though this limits us to models that run efficiently in browsers.

\textbf{Calibration Approach}: Requiring explicit calibration improves accuracy but adds friction to the user experience. Auto-calibration during initial frames was considered but proved less reliable.

\textbf{Smoothing vs. Responsiveness}: Aggressive smoothing reduces jitter but increases lag. Our sustained change detection strikes a balance, filtering noise while maintaining responsive control.

\textbf{Feature Breadth vs. Depth}: We implemented multiple facial features (head pose, mouth, eyes, eyebrows) rather than deeply optimizing single features, prioritizing versatility over perfection.

\subsection{Limitations and Challenges}

Several limitations were encountered:

\textbf{Occlusion Handling}: MediaPipe loses tracking if the face is significantly occluded or turned beyond $\sim$60° from frontal view. Hand movements near the face can cause temporary tracking failures.

\textbf{Multi-Face Scenarios}: Our implementation processes only the first detected face. Multiple users cannot control the same application simultaneously without significant architectural changes.

\textbf{Expression Subtlety}: Detecting subtle expressions remains challenging. Current thresholds favor exaggerated movements to minimize false positives.

\textbf{Accessibility Constraints}: Users with limited facial mobility may find certain controls difficult. Customizable sensitivity per detector could address this.

\textbf{Browser Compatibility}: While WebGL and getUserMedia are widely supported, some older browsers or privacy-focused configurations may block camera access or lack necessary APIs.

\section{Conclusion and Future Work}
\label{sec:conclusion}

\subsection{Summary}

This paper presented a modular architecture for real-time face tracking and 3D object control in web browsers. By separating detection logic from application-specific mappings, we created a reusable framework that enables rapid development of interactive experiences. Three diverse applications---rotation control, first-person movement, and a Pac-Man game---demonstrated the architecture's versatility.

Our implementation achieves real-time performance on consumer hardware, robust detection through calibration and filtering, and educational transparency through well-documented modular code. The system requires only a webcam and modern browser, making sophisticated face tracking accessible without specialized equipment or software installation.

\subsection{Contributions to the Field}

This work contributes to human-computer interaction research by:

\begin{enumerate}
\item \textbf{Demonstrating practical browser-based face tracking} at interactive frame rates using freely available libraries
\item \textbf{Proposing a modular architecture pattern} applicable to other computer vision applications
\item \textbf{Providing quantitative analysis} of detection accuracy and performance characteristics
\item \textbf{Creating an open-source toolkit} for developers to easily integrate face tracking into games and interactive applications
\end{enumerate}

The complete codebase serves as a reference implementation for developers exploring face tracking applications, with documentation suitable for educational contexts.

\subsection{Future Directions}

Several promising directions for future work include:

\textbf{Advanced Detectors}: Implementing detection for additional facial features such as nose wrinkle, tongue protrusion, or individual eyebrow control would expand the interaction vocabulary.

\textbf{Machine Learning Enhancement}: Training custom models to recognize user-specific expressions or gestures could improve accuracy and enable personalized interaction styles.

\textbf{Multi-User Support}: Extending the architecture to handle multiple simultaneous faces would enable collaborative applications and multi-player games.

\textbf{Adaptive Sensitivity}: Implementing automatic sensitivity adjustment based on user performance could improve accessibility and reduce calibration requirements.

\textbf{Cross-Platform Optimization}: Investigating platform-specific optimizations (e.g., WebGPU for graphics, WebAssembly for compute) could improve performance on lower-end devices.

\textbf{Application Domains}: Exploring additional application areas such as accessibility tools (computer cursor control), creative tools (3D modeling through gestures), or therapeutic applications (facial exercise games for stroke rehabilitation).

\textbf{Hybrid Input}: Combining face tracking with other input modalities (voice, hand gestures via MediaPipe Hands, traditional controllers) could create richer interaction possibilities.

\textbf{Longitudinal User Studies}: Formal user studies comparing face tracking to traditional input methods across different task types would provide empirical evidence for the approach's effectiveness.

\subsection{Broader Impact}

Face tracking technology raises important considerations:

\textbf{Privacy}: Our client-side approach ensures facial data never leaves the user's device, an important consideration given growing privacy concerns around biometric data.

\textbf{Accessibility}: Face tracking offers alternative input methods for users with limited hand mobility, though care must be taken to accommodate varying facial mobility as well.

\textbf{Inclusion}: The system's reliance on facial landmark detection may perform differently across demographic groups. Ensuring robust performance across diverse users requires careful testing and potential model refinement.

\textbf{Creative Expression}: Enabling new forms of interactive art and gaming experiences that respond to natural human expression rather than mechanical inputs.

As face tracking becomes more prevalent, maintaining transparency, respecting user privacy, and ensuring inclusive design will be essential responsibilities for developers in this space.

\subsection{Final Remarks}

The democratization of computer vision through web APIs and efficient neural networks has made sophisticated face tracking accessible to web developers and students. Our modular architecture demonstrates that production-quality interactive experiences can be built using these technologies while maintaining code clarity and reusability.

The three example applications show that the same fundamental detection capabilities can power wildly different interaction paradigms through thoughtful data mapping. This separation of concerns---what is detected versus how detections are interpreted---represents a powerful design pattern applicable beyond face tracking to other sensor-driven applications.

We hope this work inspires further exploration of natural, expressive interfaces and demonstrates that complex computer vision systems can be both powerful and comprehensible. The complete open-source implementation is available for the community to build upon, extend, and adapt to new creative and practical applications.

\section*{Acknowledgments}

AI assistance from Claude (Anthropic) and GitHub Copilot was used during the development and debugging of the codebase. All architectural decisions, algorithm implementations, and creative applications were designed and validated by the author.

\bibliographystyle{ACM-Reference-Format}
\begin{thebibliography}{22}

\bibitem{card1983psychology}
Stuart K. Card, Allen Newell, and Thomas P. Moran. 1983.
\newblock \emph{The Psychology of Human-Computer Interaction}.
\newblock Lawrence Erlbaum Associates.

\bibitem{bazarevsky2019blazeface}
Valentin Bazarevsky, Yury Kartynnik, Andrey Vakunov, Karthik Raveendran, and Matthias Grundmann. 2019.
\newblock BlazeFace: Sub-millisecond Neural Face Detection on Mobile GPUs.
\newblock \emph{arXiv preprint arXiv:1907.05047}.

\bibitem{kartynnik2019real}
Yury Kartynnik, Artsiom Ablavatski, Ivan Grishchenko, and Matthias Grundmann. 2019.
\newblock Real-time Facial Surface Geometry from Monocular Video on Mobile GPUs.
\newblock \emph{arXiv preprint arXiv:1907.06724}.

\bibitem{wobbrock2011ability}
Jacob O. Wobbrock, Shaun K. Kane, Krzysztof Z. Gajos, Susumu Harada, and Jon Froehlich. 2011.
\newblock Ability-based design: Concept, principles and examples.
\newblock \emph{ACM Transactions on Accessible Computing} 3, 3 (2011), 1--27.

\bibitem{viola2004robust}
Paul Viola and Michael J. Jones. 2004.
\newblock Robust real-time face detection.
\newblock \emph{International Journal of Computer Vision} 57, 2 (2004), 137--154.

\bibitem{sun2013deep}
Yi Sun, Xiaogang Wang, and Xiaoou Tang. 2013.
\newblock Deep convolutional network cascade for facial point detection.
\newblock In \emph{Proceedings of the IEEE Conference on Computer Vision and Pattern Recognition}. 3476--3483.

\bibitem{cao2014face}
Xudong Cao, Yichen Wei, Fang Wen, and Jian Sun. 2014.
\newblock Face alignment by explicit shape regression.
\newblock \emph{International Journal of Computer Vision} 107, 2 (2014), 177--190.

\bibitem{lugaresi2019mediapipe}
Camillo Lugaresi, Jiuqiang Tang, Hadon Nash, Chris McClanahan, Esha Uboweja, Michael Hays, Fan Zhang, Chuo-Ling Chang, Ming Guang Yong, Juhyun Lee, Wan-Teh Chang, Wei Hua, Manfred Georg, and Matthias Grundmann. 2019.
\newblock MediaPipe: A framework for building perception pipelines.
\newblock \emph{arXiv preprint arXiv:1906.08172}.

\bibitem{bolt1980put}
Richard A. Bolt. 1980.
\newblock Put-That-There: Voice and gesture at the graphics interface.
\newblock \emph{ACM SIGGRAPH Computer Graphics} 14, 3 (1980), 262--270.

\bibitem{shotton2011real}
Jamie Shotton, Andrew Fitzgibbon, Mat Cook, Toby Sharp, Mark Finocchio, Richard Moore, Alex Kipman, and Andrew Blake. 2011.
\newblock Real-time human pose recognition in parts from single depth images.
\newblock In \emph{Proceedings of the IEEE Conference on Computer Vision and Pattern Recognition}. 1297--1304.

\bibitem{majaranta2014eye}
Päivi Majaranta and Andreas Bulling. 2014.
\newblock Eye tracking and eye-based human-computer interaction.
\newblock In \emph{Advances in Physiological Computing}. 39--65.

\bibitem{tian2001recognizing}
Ying-li Tian, Takeo Kanade, and Jeffrey F. Cohn. 2001.
\newblock Recognizing action units for facial expression analysis.
\newblock \emph{IEEE Transactions on Pattern Analysis and Machine Intelligence} 23, 2 (2001), 97--115.

\bibitem{murphy2009head}
Erik Murphy-Chutorian and Mohan M. Trivedi. 2009.
\newblock Head pose estimation in computer vision: A survey.
\newblock \emph{IEEE Transactions on Pattern Analysis and Machine Intelligence} 31, 4 (2009), 607--626.

\bibitem{betke2002camera}
Margrit Betke, James Gips, and Peter Fleming. 2002.
\newblock The Camera Mouse: visual tracking of body features to provide computer access for people with severe disabilities.
\newblock \emph{IEEE Transactions on Neural Systems and Rehabilitation Engineering} 10, 1 (2002), 1--10.

\bibitem{teather2013pointing}
Robert J. Teather and Wolfgang Stuerzlinger. 2013.
\newblock Pointing at 3D target projections with one-eyed and stereo cursors.
\newblock In \emph{Proceedings of the SIGCHI Conference on Human Factors in Computing Systems}. 159--168.

\bibitem{ware1997selection}
Colin Ware and Kevin Lowther. 1997.
\newblock Selection using a one-eyed cursor in a fish tank VR environment.
\newblock \emph{ACM Transactions on Computer-Human Interaction (TOCHI)} 4, 4 (1997), 309--322.

\bibitem{zhou2022depth}
Qian Zhou, George Fitzmaurice, and Fraser Anderson. 2022.
\newblock In-Depth Mouse: Integrating Desktop Mouse into Virtual Reality.
\newblock In \emph{Proceedings of the 2022 CHI Conference on Human Factors in Computing Systems}. 1--17.

\bibitem{haas2017bringing}
Andreas Haas, Andreas Rossberg, Derek L. Schuff, Ben L. Titzer, Michael Holman, Dan Gohman, Luke Wagner, Alon Zakai, and JF Bastien. 2017.
\newblock Bringing the web up to speed with WebAssembly.
\newblock \emph{ACM SIGPLAN Notices} 52, 6 (2017), 185--200.

\bibitem{smilkov2019tensorflow}
Daniel Smilkov, Nikhil Thorat, Yannick Assogba, Ann Yuan, Nick Kreeger, Ping Yu, Kangyi Zhang, Shanqing Cai, Eric Nielsen, David Soergel, Stan Bileschi, Michael Terry, Charles Nicholson, Sandeep N. Gupta, Sarah Sirajuddin, D. Sculley, Rajat Monga, Greg Corrado, Fernanda B. Viégas, and Martin Wattenberg. 2019.
\newblock TensorFlow.js: Machine learning for the web and beyond.
\newblock \emph{arXiv preprint arXiv:1901.05350}.

\bibitem{cao2016real}
Chen Cao, Hongzhi Wu, Yanlin Weng, Tianjia Shao, and Kun Zhou. 2016.
\newblock Real-time facial animation with image-based dynamic avatars.
\newblock \emph{ACM Transactions on Graphics} 35, 4 (2016), 1--12.

\bibitem{parnas1972criteria}
David L. Parnas. 1972.
\newblock On the criteria to be used in decomposing systems into modules.
\newblock \emph{Communications of the ACM} 15, 12 (1972), 1053--1058.

\bibitem{soukupova2016real}
Tereza Soukupová and Jan Čech. 2016.
\newblock Real-time eye blink detection using facial landmarks.
\newblock In \emph{21st Computer Vision Winter Workshop}. 1--8.

\end{thebibliography}

\end{document}
